\par\indent\SourceLeaf{229}%
The \textit{second} method expresses the coordinates as functions of two variables \(p,q\). Suppose that differentiating these functions gives
\[
\begin{aligned}
\mathrm{d}x&=a\mathrm{d}p+a'\mathrm{d}q\\
\mathrm{d}y&=b\mathrm{d}p+b'\mathrm{d}q\\
\mathrm{d}z&=c\mathrm{d}p+c'\mathrm{d}q
\end{aligned}
\]
\noindent Substituting these values into the formula given above, we obtain
\[
(aX+bY+cZ)\mathrm{d}p+(a'X+b'Y+c'Z)\mathrm{d}q=0
\]
\noindent Since this equation must hold independently of the values of the differentials \(\mathrm{d}p,\mathrm{d}q\), we must evidently have
\[
aX+bY+cZ=0,\qquad a'X+b'Y+c'Z=0
\]
\noindent from which we conclude that \(X,Y,Z\) must be proportional to the quantities
\[
bc'-cb',\quad ca'-ac',\quad ab'-ba'
\]
\par\indent Thus, setting for brevity
\[
\BareRoot{((bc'-cb')^2+(ca'-ac')^2+(ab'-ba')^2)}=\Delta
\]
\noindent we have either
\[
X=\frac{bc'-cb'}{\Delta},\quad Y=\frac{ca'-ac'}{\Delta},\quad Z=\frac{ab'-ba'}{\Delta}
\]
\noindent or
\[
X=\frac{cb'-bc'}{\Delta},\quad Y=\frac{ac'-ca'}{\Delta},\quad Z=\frac{ba'-ab'}{\Delta}
\]
\par\indent Besides these two general methods there is a \textit{third}, in which one coordinate, say \(z\), is expressed as a function of the other two \(x,y\). This method is evidently a special case of either the first or the second. If we set
\[
\mathrm{d}z=t\mathrm{d}x+u\mathrm{d}y
\]
\noindent we have either
\[
X=\frac{-t}{\BareRoot{(1+tt+uu)}},\quad Y=\frac{-u}{\BareRoot{(1+tt+uu)}},\quad Z=\frac{1}{\BareRoot{(1+tt+uu)}}
\]
\noindent or
\[
X=\frac{t}{\BareRoot{(1+tt+uu)}},\quad Y=\frac{u}{\BareRoot{(1+tt+uu)}},\quad Z=\frac{-1}{\BareRoot{(1+tt+uu)}}
\]
